\section{Italian Law n.48/2008}

Italian Law n. 48/2008 represents the formal ratification of the \textbf{Budapest Convention on Cybercrime} (2001). While the Convention provides the international baseline for harmonizing cybercrime legislation, the Italian implementation focuses on the modernization of the legal system to address the rapidly evolving digital landscape.
\bigskip

\noindent
The law introduced three primary pillars of innovation into the Italian legal system:
\begin{itemize}
    \item \textbf{International Harmonization:} Alignment with the 68 member states of the Budapest Convention to facilitate cross-border cooperation.
    \item \textbf{Reorganization of Cybercrime Offenses:} Integration of specific digital crimes into the Italian Penal Code to address modern IT threats.
    \item \textbf{Corporate Liability:} Extension of \textbf{Legislative Decree 231/2001} to include cybercrime, establishing entity liability for offenses committed in the organization's interest.
\end{itemize}

\subsection{Typology of Digital Offenses}
The 2008 reform categorizes offenses based on their impact on the CIA triad.

\subsubsection*{Unauthorized Access (Article 615-ter)}
The concept of \textit{"Accesso Abusivo"} is centered on the principle of \textit{ius excludendi alios} (the right to exclude others). 
\begin{itemize}
    \item \textbf{Authorization Limits:} A crime is committed not only by bypassing security but also by remaining in a system against the owner's will.
    \item \textbf{Subjective Scope:} Even authorized users (e.g., employees) commit an offense if they access data for purposes outside their professional mandate, such as private interest or unauthorized disclosure.
\end{itemize}

\subsubsection*{Damage and Availability}
The legal definition of "damage" has shifted from physical destruction to functional unavailability.
\begin{itemize}
    \item \textbf{Temporary Unavailability:} Actions that render a system temporarily unusable (e.g., DDoS or Ransomware) constitute criminal damage.
    \item \textbf{Data Integrity:} Deleting data from corporate devices, even if backups exist in the cloud, is punishable if the intent was to create unavailability or violate company policy.
\end{itemize}

\begin{table}[H]
\centering
\small
\begin{tabular}{|p{3cm}|p{5cm}|p{2.5cm}|}
\hline
\textbf{Crime Type} & \textbf{Legal Principle / Article} & \textbf{Primary Focus} \\
\hline
Unauthorized Access & \textit{Ius excludendi alios} (Art. 615-ter) & Confidentiality \\
System Damage & Temporary/Permanent Unavailability & Availability \\
Data Alteration & Fraudulent intent or public impact & Integrity \\
Deepfakes & AI Law 2025 (Advanced Falsification) & Integrity/Trust \\
\hline
\end{tabular}
\caption{Classification of digital offenses under Law 48/2008 and 2025 updates.}
\end{table}

\noindent
A significant limitation of Law 48/2008 is its nature as a framework of principles rather than a detailed procedural manual.
\begin{itemize}
    \item \textbf{Jurisprudential Reliance:} In the absence of specific technical decrees, the application of the law relies heavily on Court of Cassation rulings.
    \item \textbf{Dynamic Interpretation:} Legal standards regarding "temporary unavailability" and "legal access" evolve annually through case law, requiring continuous compliance monitoring.
\end{itemize}


\subsection{ \textit{Ius Excludendi Alios}}
The core of Article 615-ter of the Italian Penal Code is the principle of \textit{ius excludendi alios}, the right of the system owner to exclude others. This legal framework creates a direct parallel between the protection of a physical domicile and a digital system.

\subsubsection*{Limit of Authorization}
A critical interpretation by the Italian Court of Cassation concerns users who possess legitimate credentials but use them for unauthorized purposes.
\begin{itemize}
    \item \textbf{Functional Violation:} An official (e.g., from the Tax Agency) who has technical access to a file but consults it for personal reasons (e.g., a favor for a relative) commits \textit{unauthorized access}.
    \item \textbf{Breach of Duty:} The crime occurs the moment the actor violates the specific instructions or professional duties defined by the data controller.
    \item \textbf{Formal Nature:} Unauthorized access is a formal crime; it does not require proof of digital damage or data exfiltration to be prosecuted.
\end{itemize}

\subsection{Evolution of Digital Damage and Data Alteration}
The legal definition of "damage" in the Italian system has evolved from a 1944-era focus on physical destruction to a modern understanding of functional availability.
\medskip

\noindent
The Italian judiciary now recognizes that "damage" does not require permanent destruction.
\begin{itemize}
    \item \textbf{Temporary Unavailability:} Actions such as DDoS attacks or Ransomware that render a system unusable for a period constitute criminal damage.
    \item \textbf{Intentionality:} In a criminal trial, the prosecution must demonstrate both the material element (unavailability) and the intent (\textit{mens rea}) of the defendant.
    \item \textbf{The Deletion Paradox:} Deleting personal data from a corporate device can be a crime if it violates company policy. Even if the data exists in a cloud backup, the act of local deletion with the intent to create a "tricky issue" for the employer can lead to criminal proceedings.
\end{itemize}

\subsubsection*{Alteration and Falsification}
Since 2016, the Italian system has narrowed the scope of criminal data alteration.
\begin{itemize}
    \item \textbf{Private vs. Public:} Falsifying information between private parties is generally no longer a crime unless it constitutes fraud. However, falsifying documents for public institutions remains a serious offense.
    \item \textbf{Deepfakes:} The 2025 AI legislation formally introduced specific crimes related to deepfakes to address the modern ease of digital falsification.
\end{itemize}


\begin{figure}[H]
\centering
\begin{tikzpicture}[font=\small,node distance=6mm,>=Latex,
box/.style={rectangle,rounded corners,draw=black,fill=gray!10,align=center,text width=5cm,minimum height=5mm}]
\node[box] (step1) {\textbf{Legal Access}\\Legitimate Credentials};
\node[box, below=0.5cm of step1] (step2) {\textbf{Policy Violation}\\Use for unauthorized purpose};
\node[box, below=0.5cm of step2, fill=red!5] (step3) {\textbf{Art. 615-bis/ter}\\Criminal Unauthorized Access};
\draw[->] (step1) -- (step2);
\draw[->] (step2) -- (step3);
\end{tikzpicture}
\caption{The shift from legitimate access to criminal offense based on functional intent.}
\end{figure}


\subsection{Digital Forensics under Law 48}
The ratification of the Budapest Convention through Law 48/2008 introduced rigorous procedural requirements for the identification, collection, and preservation of digital evidence to ensure its admissibility in court.
\medskip

\noindent
Digital forensics is governed by the need to maintain the chain of custody and ensure that the evidence remains unaltered from the moment of acquisition.
\begin{itemize}
    \item \textbf{Identification vs. Acquisition:} The procedure distinguishes between identifying relevant data and the technical act of "freezing" it in its original state.
    \item \textbf{Integrity through Hashing:} To guarantee immutability, every digital find must be associated with a unique hash value. Any alteration of the data during acquisition is subject to criminal sanctions.
    \item \textbf{Non-Repeatable Acts:} Since digital systems are volatile, initial inspections are often classified as non-repeatable acts, requiring specific technical measures to prevent data degradation.
\end{itemize}

\subsubsection{Procedural Innovations: Inspections and Searches}
The law updated the Italian Code of Criminal Procedure to address the peculiarities of "virtual" spaces.
\begin{itemize}
    \item \textbf{Remote Inspections:} Law enforcement can perform inspections of computer systems even remotely, provided they use techniques that ensure the original data is not modified.
    \item \textbf{Duty of Cooperation:} System administrators and service providers are legally obligated to assist authorities in accessing encrypted or protected data.
\end{itemize}

\subsubsection{Data Retention}
The legal framework imposes specific obligations on electronic communication providers regarding the storage of metadata for investigative purposes.

\subsubsection*{Traffic Data Retention}
Providers must store logs to allow the retrospective reconstruction of digital crimes.
\begin{itemize}
    \item \textbf{Metadata focus:} Retention applies to the "envelope" of the communication (who, when, where) rather than the content itself.
    \item \textbf{Access Limits:} Access to this data requires a judicial mandate, balancing investigative needs with the right to privacy as outlined in the GDPR and NIS2 frameworks.
\end{itemize}

\begin{table}[H]
\centering
\small
\begin{tabular}{|p{3cm}|p{5cm}|p{2.5cm}|}
\hline
\textbf{Forensic Phase} & \textbf{Technical Requirement} & \textbf{Legal Goal} \\
\hline
Identification & System mapping and volatile data triage & Completeness \\
Acquisition & Bit-stream imaging and Hashing & Immutability \\
Preservation & Controlled storage and Chain of Custody & Admissibility \\
Analysis & Use of validated forensic tools & Reliability \\
\hline
\end{tabular}
\caption{The standard forensic workflow under Law 48/2008 requirements.}
\end{table}

\begin{figure}[H]
\centering
\begin{tikzpicture}[font=\small,node distance=7mm,>=Latex,
box/.style={rectangle,rounded corners,draw=black,fill=gray!10,align=center,text width=5cm,minimum height=6mm}]
\node[box] (step1) {\textbf{Digital Scene}\\Identification of systems};
\node[box, below=of step1] (step2) {\textbf{Acquisition}\\Hashing and Bit-copy};
\node[box, below=of step2] (step3) {\textbf{Legal Proof}\\Court Admissibility};

\draw[->] (step1) -- (step2);
\draw[->] (step2) -- node[right]{Chain of Custody} (step3);
\end{tikzpicture}
\caption{Linear flow of digital evidence from acquisition to court admission.}
\end{figure}

\subsubsection*{The Role of the CISO in Investigations}
In the context of corporate liability (D.Lgs. 231/2001), the CISO must ensure that internal incident response procedures do not contaminate potential evidence. Improper handling of a system after a breach can lead to the inadmissibility of evidence, hindering the organization's ability to defend itself or prosecute the attacker.

